
\section{Conclusion}\label{sec:conclusion}

This paper reports the results of more than 1,000 auction experiments with LLM agents. 
In particular, we find behavior that conforms with important experimental results from labs involving human
subjects (e.g., evidence of risk-averse bidding and evidence that clock auctions are easier to play). 
The results are encouraging and we see the main contribution of this work as 
 putting forward a framework for how to think about LLM experimental agents as a proxy for human agents. In particular, the design space for prompting is large, and we hope that future research 
 will make  use of our code to run simulations testing their own prompt variations. 

While this paper focuses on auction design, future work may use LLM sandboxes to test other kinds of
economic mechanisms; e.g., voting, matching, contract design, etc.
%
As  LLM models are increasingly validated as
proxies for human behavior, LLM agents can then be used to obtain what would otherwise be prohibitively expensive evidence. 
As a provocative example, while ethical and financial constraints make it impossible to run voting experiments at the scale of nations, it may be possible to run such experiments with LLM agents. 

We are particularly interested in the use of LLM-based techniques to generate synthetic data that is useful for informing 
economic design. 
In particular, some auction formats, such as combinatorial auctions, are complex and  difficult to run frequently and at scale in traditional lab experiments. 
Augmenting these traditional lab experiments with LLM-agent  experiments, when correctly validated, may open up wide new avenues to better understand the design tradeoffs in these kinds of complex and often high-stakes environments. 
